%\myChapterStar{Titre}{Titre court}{Ajouter à la table des matières? (false|true|chapter|section|subsection|subsubsection -chapter par défaut-)}
\myChapterStar{Introduction}{}{true}
%\myMinitoc{Profondeur de la minitoc (section|subsection|subsubsection)}{Titre de la minitoc}
\myMiniToc{}{Sommaire}

% \newpage
%\mySectionStar{Titre}{Titre court}{Ajouter à la table des matières? (false|true|chapter|section|subsection|subsubsection -section par défaut-)}
\mySectionStar{Contexte}{}{true}
Le paysage logiciel a considérablement évolué, passant d’applications monolithiques à une approche de plus en plus modulaire. Les applications modernes sont construites à partir de composants indépendants et agnostiques, appelés conteneurs \cite{merkel2014docker}. Contrairement aux composants traditionnels, les conteneurs encapsulent tout le code source ainsi que toutes les bibliothèques nécessaires à leur exécution. Cela permet un déploiement simplifié dans tous les environnements, à la fois locaux et cloud. Le déploiement d’applications conteneurisées se fait au travers du processus d’orchestration. L’orchestration est le processus automatisé de planification et de gestion des déploiements de conteneurs dans un cluster de machines. Il tient compte de diverses exigences de service, telles que la mémoire, la puissance de calcul (nombre de processeurs), les répliques (copies) souhaitées ou même l’emplacement physique de l’instance de conteneur. L’orchestration a été largement adoptée par les entreprises technologiques mainstream telles que Google, AWS et IBM\cite{app12125793,awada2018application}. En plus de tirer parti de l’orchestration de conteneurs pour leurs propres applications, elles proposent également à d’autres entreprises des services d’orchestration basés sur le cloud.
Un orchestrateur particulier, connu sous le nom de Kubernetes\cite{7036275}, s’est largement répandu. Développé conjointement par Google et la Fondation Linux en 2014, Kubernetes utilise le concept de pod comme unité fondamentale de calcul au sein d’un cluster. Un pod peut héberger plusieurs conteneurs, et un service est déployé en répliquant les pods à travers le cluster, chacun contenant les conteneurs nécessaires à l’exécution du service. Cela présente plusieurs avantages. Fiabilité accrue : en configurant plus de répliques dans les spécifications de déploiement de Kubernetes, un service devient plus résilient aux pannes. Déploiement continu simplifié : les mises à jour et les retours en arrière pour les services conteneurisés sont optimisés de façon à avoir des temps d’arrêt de
services nuls.

Une innovation récente de l’équipe Fuchsia en 2024 \cite{ngoufo2024} introduit le concept de composition incrémentale des services au sein de Kubernetes. Cela permet de définir de nouveaux services en combinant paresseusement des services existants déjà déployés dans le cluster. Ces services composites sont essentiellement des ensembles de coroutines (processus conçus pour le travail parallèle et coopératif). Chaque coroutine peut être soit un service ordinaire (service Kubernetes déjà déployé ou programme Java compilé), soit un autre service composite lui-même. Le principal avantage d’une telle composition de services réside dans sa nature incrémentale. Les services composites traitent leurs entrées de manière paresseuse et fournissent leurs résultats de manière incrémentale. Cette efficacité est obtenue en tirant parti de la terminaison individuelle des coroutines. Chaque coroutine fonctionne de manière indépendante et dirigée par les données. Lorsque certaines de ses données d’entrée deviennent disponibles, la coroutine s’active, produisant les sorties en fonction des entrées disponibles. Le service composite peut alors s’appuyer sur les sorties produites pour renvoyer certaines de ses propres sorties. L’équipe Fuchsia a utilisé ce concept de calcul incrémental pour déployer un service composite de tri sur un cluster virtuel d’un ordinateur portable moyen de gamme. Les résultats ont montré que la composition incrémentale de services sur Kubernetes présente de bonnes performances en matière de distribution, de parallélisme et de robustesse. À cette fin, l’équipe a développé un langage dédié à la composition de services sur Kubernetes. Le langage développé combine la syntaxe Yaml au langage de programmation Java. Les résultats présentés par l’équipe Fuchsia sont très prometteurs. Ils offrent la possibilité de composer et de distribuer des services de la même manière que l’on écrirait un programme dans un langage générique, avec la différence que les procédures sont remplacées par des services.


\mySectionStar{Problématique}{}{true}
% Il est crucial de fournir une guidance structurée pour aider les nouveaux étudiants à naviguer efficacement dans l'abondance des formations offertes par les universités.
% La problématique découlante est celle de savoir  Comment concevoir une démarche formelle pouvant contribuer à la création d’un modèle d’apprentissage automatique efficace pour la recommandation des formations universitaires ?
Dans les architectures logicielles modernes, la composition de services permet de construire des processus complexes à partir d’unités logicielles autonomes et réutilisables. Les approches classiques de composition, qu’elles soient statiques ou dynamiques, reposent généralement sur une orchestration centralisée et une définition globale du processus, ce qui limite leur capacité à s’adapter à des environnements fortement distribués, évolutifs et pilotés par les données.

Or, les infrastructures cloud-native actuelles — \textit{où les services sont conteneurisés, distribués sur plusieurs nœuds, et déclenchés en fonction d’événements asynchrones} — requièrent des modèles de composition plus flexibles, capables de réagir localement aux changements de contexte, de s’exécuter partiellement et de manière paresseuse, et de tirer parti du parallélisme inhérent à certaines applications.

La composition incrémentale des services, fondée sur les grammaires attribuées gardées (GAG), constitue une réponse prometteuse à ces enjeux. Cependant, cette approche reste encore peu étudiée d’un point de vue expérimental : ses performances, sa robustesse, et sa capacité à exploiter efficacement les ressources cloud-native doivent être évaluées concrètement à travers des cas d’usage pertinents.

C’est dans cette perspective que s’inscrit la problématique de notre travail qui est de savoir comment concevoir, déployer et mesurer l’efficacité d’un moteur de composition incrémentale dans un environnement distribué?


\mySectionStar{Objectif}{}{true}
 % \textbf{Objectif principal} :
  L’objectif de ce travail consiste à mettre en place et valider un benchmark reproductible de la transformée de Fourier, afin de comparer de manière rigoureuse les performances des compositions de services incrémentale et non incrémentale dans un environnement distribué. \\
  % L’objectif principal de ce travail consiste à mettre en place et valider un benchmark reproductible de la transformée de Fourier, afin de comparer de manière rigoureuse les performances des compositions de services incrémentale et non incrémentale dans un environnement distribué. \\
  % La réalisation de cet objectif passe par plusieurs sous objectifs ou objectifs spécifiques. \\
  
%   \textbf{Objectifs spécifiques:}
% \begin{enumerate}
%     \item \textbf{Développer un modèle mathématique} :
%     \begin{itemize}
%         \item Élaborer les formules mathématiques pour décrire les opérations de base.
%         \item Définir la métrique le calcul des distances.
%     \end{itemize}
% \end{enumerate}
    

\mySectionStar{Structure du document}{}{true}
La suite de ce document se compose de trois chapitres et d'une conclusion générale, qui abordent les points suivants de manière approfondies :\\
\\
\textbf{Chapitre I : Généralités sur la composition des services}: Dans ce chapitre, nous présenterons les aspects importants de l'architecture orientée services et des microservices puis nous présenterons une vue approfondie de l'algorithme de transformée de Fourier.
\\
\\
\textbf{Chapitre II : État de l’art sur la composition incrémentale des services}: Nous y explorerons les limites des approches classiques de composition, puis nous introduisons les principes de la composition incrémentale. Une attention particulière est portée au modèle formel des Grammaires Attribuées Gardées (GAG), ainsi qu’aux travaux de recherche qui ont contribué à son développement.
\\
\\
\textbf{Chapitre III : Contribution au benchmark d’évaluation de la composition incrémentale}: Nous présenterons spécifiquement la démarche entreprise depuis l'approche GAG de l'algorithme jusqu'au déploiement sur un cluster Kubernetes et l'analyse des résultats.
\\
\\
\textbf{Conclusion générale :} Le bilan du travail se fera par le rappel de la problématique, une analyse critique des résultats obtenus, puis la présentation des perspectives.

\myCleanStarChapterEnd
